\documentclass[11pt]{article}

\usepackage[a4paper,margin=1in]{geometry}
\usepackage{amsmath,amssymb,amsthm,bm}
\usepackage{booktabs}
\usepackage{array}
\usepackage{tabularx}
\usepackage{enumitem}
\usepackage{xcolor}
\usepackage{hyperref}
\usepackage{microtype}
\usepackage{xspace}

\hypersetup{
    colorlinks=true,
    linkcolor=blue!50!black,
    citecolor=blue!50!black,
    urlcolor=blue!50!black
}

\newtheorem{assumption}{Assumption}
\newtheorem{remark}{Remark}
\newtheorem{definition}{Definition}

\newcommand{\A}{\mathcal A}
\newcommand{\F}{\mathcal F}
\newcommand{\Lcal}{\mathcal L}
\newcommand{\R}{\mathbb R}
\newcommand{\E}{\mathbb E}
\newcommand{\Prob}{\mathbb P}
\newcommand{\norm}[1]{\left\lVert #1 \right\rVert}
\newcommand{\argmax}{\operatorname*{arg\,max}}
\newcommand{\argmin}{\operatorname*{arg\,min}}
\newcommand{\LCB}{\operatorname{LCB}}
\newcommand{\UCB}{\operatorname{UCB}}
\newcommand{\ID}{\textsc{Identify}\xspace}
\newcommand{\CERT}{\textsc{Certify}\xspace}
\newcommand{\COMMIT}{\textsc{Commit}\xspace}
\newcommand{\TS}{\mathrm{TS}}

\title{Library Identify-and-Commit with Thompson-Guided Certification\\\large Detailed Algorithmic Specification}
\author{}
\date{}

\begin{document}
\maketitle

\begin{abstract}
This note defines a recurrent-regime bandit algorithm for settings in which a finite, initially unknown, collection of operating regimes is encountered repeatedly. Each regime is represented by a centered reward shape, which removes common reward-level shifts. The algorithm first identifies the active regime from balanced blocks. If the regime has a previously certified arm, it commits to that arm immediately. If the best arm is still uncertain, it enters a certification phase. The certification phase uses a top-two Thompson sampling rule only to decide which arm to sample next, and it stops according to a frequentist $\varepsilon_{\mathrm{opt}}$-optimality condition. Once an arm is certified as practically optimal, the algorithm commits to it and monitors its reward for drift. Thus Thompson sampling is not used in stable operation; it is used only as an adaptive sampling rule for certification.
\end{abstract}

\section{Setting}

There are $K$ arms,
\begin{equation}
    \A=\{1,\ldots,K\}.
\end{equation}
At time $t$, the learner selects an arm $A_t\in\A$ and observes a reward $Y_t\in\R$. The environment is piecewise stationary. During a stable episode, rewards are generated by an unobserved regime $z$. Regimes may recur: after the system leaves regime $z$, it may later return to the same regime.

For a stable regime $z$, the arm means are written as
\begin{equation}
    \mu_k^{(z)}=\ell^{(z)}+\theta_k^{(z)},
    \qquad k\in\A,
    \label{eq:level-shape}
\end{equation}
where $\ell^{(z)}$ is a common reward level and $\bm\theta^{(z)}$ is the action-relevant shape. The shape is normalized by
\begin{equation}
    \sum_{k=1}^K\theta_k^{(z)}=0.
\end{equation}
The best arm in regime $z$ is therefore
\begin{equation}
    k^\star(z)\in\argmax_{k\in\A}\theta_k^{(z)}.
\end{equation}
A common upward or downward shift of all rewards changes $\ell^{(z)}$ but does not change the best arm.

For any vector $\bm v\in\R^K$, define the centering operator
\begin{equation}
    P=I-\frac{1}{K}\bm 1\bm 1^\top,
    \qquad
    \phi(\bm v)=P\bm v.
    \label{eq:centering}
\end{equation}
The algorithm identifies regimes by centered reward shapes.

\begin{assumption}[Local sub-Gaussian rewards]
\label{ass:subgaussian}
During a stable episode with active regime $Z_t$, rewards satisfy
\begin{equation}
    Y_t=\mu_{A_t}^{(Z_t)}+\varepsilon_t,
\end{equation}
where
\begin{equation}
    \E[\varepsilon_t\mid \F_{t-1},A_t,Z_t]=0,
\end{equation}
and, for all $\lambda\in\R$,
\begin{equation}
    \E\!\left[\exp(\lambda\varepsilon_t)\mid \F_{t-1},A_t,Z_t\right]
    \leq
    \exp\!\left(\frac{\lambda^2\sigma^2}{2}\right).
    \label{eq:subgaussian}
\end{equation}
\end{assumption}

The assumption is local to a regime-arm pair. Different regimes may have different means. The concentration condition is imposed only on the centered fluctuations around the corresponding regime-arm mean.

\section{Inputs and maintained quantities}

The algorithm has the following inputs:
\begin{itemize}[leftmargin=1.5em]
    \item $K$: number of arms;
    \item $m$: number of samples per arm in one balanced identification block;
    \item $\sigma$: sub-Gaussian noise proxy;
    \item $\delta\in(0,1)$: confidence level;
    \item $\varepsilon_{\mathrm{opt}}\geq 0$: practical optimality tolerance;
    \item $\nu,h>0$: tolerance and threshold of the committed-arm drift detector.
\end{itemize}
The only additional quantity relative to exact best-arm certification is $\varepsilon_{\mathrm{opt}}$. It has a direct interpretation: reward differences below $\varepsilon_{\mathrm{opt}}$ are not worth distinguishing. If $\varepsilon_{\mathrm{opt}}=0$, the algorithm returns to exact best-arm certification.

The algorithm maintains a library
\begin{equation}
    \Lcal_t=\{1,\ldots,J_t\}.
\end{equation}
Each library entry $j$ stores two types of information.

\paragraph{Shape information.}
The shape information is used for regime matching. It is updated only from balanced blocks. The entry stores
\begin{equation}
    N_j^{\mathrm{sh}},
    \qquad
    S_{j,1}^{\mathrm{sh}},\ldots,S_{j,K}^{\mathrm{sh}},
    \qquad
    \widehat{\bm\theta}^{(j)}.
\end{equation}
Here $N_j^{\mathrm{sh}}$ is the number of accepted balanced-block samples per arm assigned to regime $j$, and $S_{j,k}^{\mathrm{sh}}$ is the corresponding raw reward sum for arm $k$. The empirical vector is
\begin{equation}
    \widehat\mu_{j,k}^{\mathrm{sh}}
    =
    \frac{S_{j,k}^{\mathrm{sh}}}{N_j^{\mathrm{sh}}},
    \qquad k\in\A,
\end{equation}
and the stored centered shape is
\begin{equation}
    \widehat{\bm\theta}^{(j)}
    =
    P\widehat{\bm\mu}_{j}^{\mathrm{sh}}.
    \label{eq:library-shape}
\end{equation}

\paragraph{Certification information.}
The certification information is used to decide which arm should be played in regime $j$. It is stored in centered coordinates. The entry stores
\begin{equation}
    N_{j,k}^{\mathrm{cert}},
    \qquad
    S_{j,k}^{\mathrm{cert}},
    \qquad k\in\A,
\end{equation}
where $N_{j,k}^{\mathrm{cert}}$ is the amount of centered evidence for arm $k$ and $S_{j,k}^{\mathrm{cert}}$ is the corresponding centered reward sum. The certification estimate is
\begin{equation}
    \widehat\theta_{j,k}^{\mathrm{cert}}
    =
    \frac{S_{j,k}^{\mathrm{cert}}}{N_{j,k}^{\mathrm{cert}}}.
    \label{eq:cert-estimate}
\end{equation}
The certified arm is denoted by
\begin{equation}
    a_j\in\A\cup\{\varnothing\}.
\end{equation}
The value $a_j=\varnothing$ means that the regime is known but no arm has yet been certified as practically optimal.

Balanced blocks assigned to regime $j$ update both the shape information and the certification information. During certification, additional arm samples update only the certification information, after centering by the local level of the current identified episode. They do not update the full shape representation, because unbalanced samples do not identify the full centered reward vector under possible level shifts.

\section{Confidence radii}

For $n\geq1$, define the coordinate confidence radius
\begin{equation}
    r(n)=\sigma\sqrt{\frac{2\log(8KT/\delta)}{n}}.
    \label{eq:coord-radius}
\end{equation}
For centered shape coordinates, use
\begin{equation}
    b(n)=2r(n).
    \label{eq:shape-radius}
\end{equation}
The constants are not essential. The important point is that $b(n)$ decreases as more balanced samples are assigned to the same shape.

For certification, define arm-specific bounds
\begin{equation}
    B_{j,k}^{\mathrm{cert}}
    =
    \sigma\sqrt{\frac{2\log(8KT/\delta)}{N_{j,k}^{\mathrm{cert}}}},
    \label{eq:cert-radius}
\end{equation}
with
\begin{equation}
    \LCB_{j,k}
    =
    \widehat\theta_{j,k}^{\mathrm{cert}}-B_{j,k}^{\mathrm{cert}},
    \qquad
    \UCB_{j,k}
    =
    \widehat\theta_{j,k}^{\mathrm{cert}}+B_{j,k}^{\mathrm{cert}}.
    \label{eq:cert-bounds}
\end{equation}
These intervals are used to stop certification. They are not used as a UCB arm-selection rule.

\section{Balanced blocks and local level}

\begin{definition}[Balanced block]
A balanced block samples every arm exactly $m$ times. Its empirical reward vector is
\begin{equation}
    \widehat v_k
    =
    \frac{1}{m}\sum_{s\in\mathcal B:A_s=k}Y_s,
    \qquad k\in\A,
    \label{eq:block-vector}
\end{equation}
its local level estimate is
\begin{equation}
    \widehat\ell_{\mathcal B}
    =
    \frac{1}{K}\sum_{k=1}^K\widehat v_k,
    \label{eq:block-level}
\end{equation}
and its centered shape is
\begin{equation}
    \widehat{\bm\psi}_{\mathcal B}=P\widehat{\bm v}.
    \label{eq:block-shape}
\end{equation}
\end{definition}

Balanced blocks are used in the identification phase. They provide both a shape estimate and a local level estimate. The level estimate is later used to center certification samples from the same stable episode. If an observation $Y_t$ is collected during certification after an accepted block with level $\widehat\ell$, its centered certification reward is
\begin{equation}
    X_t=Y_t-\widehat\ell.
    \label{eq:centered-cert-reward}
\end{equation}
The variable $X_t$ is used to update the certification statistics of the selected arm.

The algorithm maintains a temporary stable aggregate $C$ during identification. The aggregate contains $n_C$ samples per arm, reward sums $S_k^C$, empirical means $\widehat\mu_k^C=S_k^C/n_C$, local level
\begin{equation}
    \widehat\ell_C=\frac{1}{K}\sum_{k=1}^K\widehat\mu_k^C,
\end{equation}
and centered shape
\begin{equation}
    \widehat{\bm\psi}^C=P\widehat{\bm\mu}^C.
\end{equation}
A new balanced block $B$ is compatible with the current aggregate $C$ if
\begin{equation}
    \norm{\widehat{\bm\psi}^{B}-\widehat{\bm\psi}^{C}}_\infty
    \leq
    b(m)+b(n_C).
    \label{eq:stable-block}
\end{equation}
If the condition holds, $B$ is added to $C$. If it fails, the previous aggregate is discarded and a new aggregate is initialized from $B$. This rule prevents data from different stable shapes from being averaged together.

\section{Regime identification}

In mode \ID, the algorithm repeatedly collects balanced blocks and maintains the stable aggregate $C$. For every library entry $j$, define
\begin{equation}
    D_j(C)=\norm{\widehat{\bm\psi}^{C}-\widehat{\bm\theta}^{(j)}}_\infty.
    \label{eq:distance}
\end{equation}
The aggregate $C$ is compatible with regime $j$ if
\begin{equation}
    D_j(C)\leq b(n_C)+b(N_j^{\mathrm{sh}}).
    \label{eq:compatible-regime}
\end{equation}
The match set is
\begin{equation}
    \mathcal M(C)=
    \left\{j\in\Lcal_t:
    D_j(C)\leq b(n_C)+b(N_j^{\mathrm{sh}})
    \right\}.
    \label{eq:match-set}
\end{equation}

\paragraph{Unique match.}
If $\mathcal M(C)=\{j\}$, the current regime is identified as $j$. The aggregate $C$ is assigned to library entry $j$. The shape statistics are updated by
\begin{equation}
    S_{j,k}^{\mathrm{sh}}
    \leftarrow
    S_{j,k}^{\mathrm{sh}}+S_k^C,
    \qquad
    N_j^{\mathrm{sh}}
    \leftarrow
    N_j^{\mathrm{sh}}+n_C,
    \qquad k\in\A.
    \label{eq:shape-update}
\end{equation}
The certification statistics are updated with the centered aggregate:
\begin{equation}
    S_{j,k}^{\mathrm{cert}}
    \leftarrow
    S_{j,k}^{\mathrm{cert}}+n_C\widehat\psi_k^C,
    \qquad
    N_{j,k}^{\mathrm{cert}}
    \leftarrow
    N_{j,k}^{\mathrm{cert}}+n_C,
    \qquad k\in\A.
    \label{eq:cert-update-from-block}
\end{equation}
The estimates and confidence bounds are recomputed. If $a_j$ is already certified and still satisfies the $\varepsilon_{\mathrm{opt}}$ rule in~\eqref{eq:epsilon-cert}, the algorithm enters \COMMIT. Otherwise it enters \CERT.

\paragraph{Multiple matches.}
If $|\mathcal M(C)|>1$, the current data do not distinguish the matched regimes. If all matched regimes have the same certified arm, the ambiguity is irrelevant for action and the algorithm may commit to that common arm. If the matched regimes have different certified arms, or if at least one matched regime has no certified arm, the algorithm remains in \ID and collects more balanced blocks.

\paragraph{No match.}
If $\mathcal M(C)=\varnothing$, the aggregate is treated as a candidate new regime. A new library entry is not created immediately. The algorithm enters \CERT with candidate $C$. The candidate becomes a stored regime only if its shape remains separated from all existing entries and an arm is certified as $\varepsilon_{\mathrm{opt}}$-optimal.

\section{Thompson-guided certification}

The certification phase has two roles. First, it decides which arm should be used for a regime. Second, it prevents the algorithm from spending excessive samples trying to distinguish arms whose means are practically tied. The phase uses Thompson sampling to choose which arm to sample next, but it stops according to a confidence rule.

\subsection{Posterior used for sampling}

Consider either an existing regime $j$ or a candidate regime $C$. In both cases, certification maintains arm-wise centered statistics
\begin{equation}
    N_k^{\mathrm{cert}},
    \qquad
    S_k^{\mathrm{cert}},
    \qquad k\in\A.
\end{equation}
For an existing regime $j$, these are $N_{j,k}^{\mathrm{cert}}$ and $S_{j,k}^{\mathrm{cert}}$. For a candidate $C$, they are initialized by
\begin{equation}
    N_{C,k}^{\mathrm{cert}}=n_C,
    \qquad
    S_{C,k}^{\mathrm{cert}}=n_C\widehat\psi_k^C,
    \qquad k\in\A.
    \label{eq:candidate-cert-init}
\end{equation}
The current certification estimate is
\begin{equation}
    \widehat\theta_k^{\mathrm{cert}}
    =
    \frac{S_k^{\mathrm{cert}}}{N_k^{\mathrm{cert}}}.
    \label{eq:local-cert-estimate}
\end{equation}
For sampling purposes, the algorithm uses the Gaussian working posterior
\begin{equation}
    \vartheta_k
    \mid \mathcal D
    \sim
    \mathcal N\!\left(
    \widehat\theta_k^{\mathrm{cert}},
    \frac{\sigma^2}{N_k^{\mathrm{cert}}}
    \right),
    \qquad k\in\A.
    \label{eq:working-posterior}
\end{equation}
This posterior is used only to choose the next arm in the certification phase. The certification guarantee comes from the confidence bounds, not from assuming that the posterior is exactly correct.

At each certification round, draw independently
\begin{equation}
    \widetilde\vartheta_k
    \sim
    \mathcal N\!\left(
    \widehat\theta_k^{\mathrm{cert}},
    \frac{\sigma^2}{N_k^{\mathrm{cert}}}
    \right),
    \qquad k\in\A,
\end{equation}
and let the \emph{leader} be
\begin{equation}
    I_t\in\argmax_{k\in\A}\widetilde\vartheta_k.
    \label{eq:ts-leader}
\end{equation}
The played arm is chosen by a top-two rule with parameter $\beta\in(0,1]$. With probability $\beta$, the leader is played, $A_t=I_t$. With probability $1-\beta$, independent posterior redraws $\widetilde\vartheta'_k$ are taken until their $\argmax$ differs from $I_t$, and the resulting \emph{challenger}
\begin{equation}
    J_t\in\argmax_{k\neq I_t}\widetilde\vartheta'_k
\end{equation}
is played, $A_t=J_t$:
\begin{equation}
    A_t=
    \begin{cases}
        I_t & \text{with probability }\beta,\\
        J_t & \text{with probability }1-\beta.
    \end{cases}
    \label{eq:ts-cert-action}
\end{equation}
The choice $\beta=1$ recovers plain Thompson sampling (the leader is always played). A value $\beta\in(0,1)$ guarantees that a fixed fraction of certification pulls is spent on a challenger arm; this is what drives down $\max_{k\neq\widehat a}\UCB_k$ and lets the stopping rule~\eqref{eq:epsilon-cert} fire (see Remark~\ref{rem:toptwo}). Unless stated otherwise we use $\beta=\tfrac12$.

After observing $Y_t$, center it using the local level $\widehat\ell_C$ of the current identified aggregate:
\begin{equation}
    X_t=Y_t-\widehat\ell_C.
\end{equation}
Then update only the selected arm:
\begin{equation}
    S_{A_t}^{\mathrm{cert}}
    \leftarrow
    S_{A_t}^{\mathrm{cert}}+X_t,
    \qquad
    N_{A_t}^{\mathrm{cert}}
    \leftarrow
    N_{A_t}^{\mathrm{cert}}+1.
    \label{eq:cert-online-update}
\end{equation}

\subsection{$\varepsilon_{\mathrm{opt}}$-optimal certification}

Let
\begin{equation}
    \widehat a
    \in
    \argmax_{k\in\A}\widehat\theta_k^{\mathrm{cert}}.
\end{equation}
Define the certification bounds
\begin{equation}
    \LCB_k
    =
    \widehat\theta_k^{\mathrm{cert}}
    -
    \sigma\sqrt{\frac{2\log(8KT/\delta)}{N_k^{\mathrm{cert}}}},
\end{equation}
\begin{equation}
    \UCB_k
    =
    \widehat\theta_k^{\mathrm{cert}}
    +
    \sigma\sqrt{\frac{2\log(8KT/\delta)}{N_k^{\mathrm{cert}}}}.
\end{equation}
The arm $\widehat a$ is certified as practically optimal if
\begin{equation}
    \LCB_{\widehat a}
    +
    \varepsilon_{\mathrm{opt}}
    \geq
    \max_{k\neq \widehat a}\UCB_k.
    \label{eq:epsilon-cert}
\end{equation}
Condition~\eqref{eq:epsilon-cert} means that, on the confidence event, the committed arm is at most $\varepsilon_{\mathrm{opt}}$ worse than the unknown best arm. If $\varepsilon_{\mathrm{opt}}=0$, the rule requires exact best-arm separation.

This rule is the main modification. When two arms are extremely close, exact certification may require a very large number of samples. The tolerance $\varepsilon_{\mathrm{opt}}$ prevents the algorithm from spending many samples to resolve a difference that is smaller than the practical resolution of the task.

\subsection{Existing regimes}

For an existing regime $j$, certification proceeds with the statistics of entry $j$. If condition~\eqref{eq:epsilon-cert} holds, the algorithm stores
\begin{equation}
    a_j=\widehat a
\end{equation}
and enters \COMMIT. If the regime already had a certified arm and the same condition still holds for that arm, the algorithm commits immediately and does not run additional Thompson-sampling certification.

If condition~\eqref{eq:epsilon-cert} does not hold, the algorithm keeps sampling according to~\eqref{eq:ts-cert-action}. The arm samples are used only for certification statistics. They are not used to update the regime-matching shape $\widehat{\bm\theta}^{(j)}$.

\subsection{Candidate new regimes}

For a candidate new regime $C$, certification uses the initialized statistics in~\eqref{eq:candidate-cert-init}. The candidate may become a new library entry only if two conditions hold. First, an arm is $\varepsilon_{\mathrm{opt}}$-certified by~\eqref{eq:epsilon-cert}. Second, the candidate shape remains separated from the existing library:
\begin{equation}
    D_j(C)>b(n_C)+b(N_j^{\mathrm{sh}})
    \qquad
    \text{for every }j\in\Lcal_t.
    \label{eq:candidate-separated}
\end{equation}
When both conditions hold, the algorithm creates a new regime $J_t+1$ with shape statistics initialized from $C$, certification statistics initialized from the candidate certification data, and certified arm $a_{J_t+1}=\widehat a$.

A candidate with nearly flat centered shape, $P\widehat{\bm v}\approx\bm 0$, is therefore treated according to the same rule as every other candidate. If all arms are indistinguishable up to $\varepsilon_{\mathrm{opt}}$, the algorithm may certify any practically optimal arm. This is not harmful for regret, because all certified arms are within the prescribed tolerance. If $\varepsilon_{\mathrm{opt}}=0$, such a candidate will not be certified unless one arm is statistically separated.

\section{Commitment and drift detection}

When the current regime $j$ has certified arm $a_j$, the algorithm enters \COMMIT. In this mode the arm-selection rule is deterministic:
\begin{equation}
    A_t=a_j.
    \label{eq:commit}
\end{equation}
There is no UCB bonus, no Thompson sampling, and no scheduled exploration during commitment.

At the beginning of the committed phase, the algorithm stores a local baseline. Let $C_{\mathrm{cur}}$ be the most recent accepted aggregate that led to the current identification. Its local level is $\widehat\ell_{C_{\mathrm{cur}}}$. The predicted reward of the committed arm is
\begin{equation}
    \widehat m_{j,a_j}^{\mathrm{cur}}
    =
    \widehat\ell_{C_{\mathrm{cur}}}
    +
    \widehat\theta_{j,a_j}^{\mathrm{cert}}.
    \label{eq:local-baseline}
\end{equation}
At committed time $t$, define
\begin{equation}
    e_t=Y_t-\widehat m_{j,a_j}^{\mathrm{cur}}.
\end{equation}
The detector statistic is
\begin{equation}
    G_t=\max\{0,G_{t-1}+|e_t|-\nu\},
    \qquad G_{t_0}=0,
    \label{eq:cusum}
\end{equation}
where $t_0$ is the start of the committed phase. A drift alarm is raised when
\begin{equation}
    G_t\geq h.
    \label{eq:alarm}
\end{equation}
After an alarm, commitment stops, $G_t$ is reset, and the algorithm returns to \ID. Post-alarm observations are not assigned to the previous regime unless they are later accepted through balanced-block identification.

\begin{assumption}[Committed-arm visibility]
\label{ass:visibility}
Whenever the environment leaves a regime $j$ while the algorithm is committed to $a_j$, the mean reward of the committed arm changes by at least
\begin{equation}
    \left|\mu_{a_j}^{(j)}-\mu_{a_j}^{(z')}\right|
    \geq
    \eta_{\min}>0,
    \label{eq:visibility}
\end{equation}
where $z'$ is the new regime.
\end{assumption}

Assumption~\ref{ass:visibility} is the price of removing periodic full-shape checks. If a regime change affects only unplayed arms and leaves the committed arm statistically unchanged, no committed-arm detector can reliably detect it.

\section{Complete algorithm}

\begin{table}[t]
\centering
\small
\renewcommand{\arraystretch}{1.25}
\begin{tabularx}{\textwidth}{>{\raggedright\arraybackslash}p{0.16\textwidth} >{\raggedright\arraybackslash}X >{\raggedright\arraybackslash}p{0.18\textwidth}}
\toprule
\textbf{Mode} & \textbf{Action and decision rule} & \textbf{Next mode} \\
\midrule
\ID & Collect balanced blocks and maintain a stable aggregate $C$. Compute the match set $\mathcal M(C)$ from~\eqref{eq:match-set}. If there is a unique match $j$, update $j$ with the centered block evidence. If there is no match, treat $C$ as a candidate new regime. If there are multiple matches, continue unless all matched regimes have the same certified arm. & \CERT or \COMMIT \\
\midrule
\CERT & Use the top-two Thompson rule to choose certification samples, according to~\eqref{eq:ts-cert-action}. Update arm-wise centered certification statistics using~\eqref{eq:cert-online-update}. Stop when the $\varepsilon_{\mathrm{opt}}$ condition~\eqref{eq:epsilon-cert} holds. & \COMMIT if certified; otherwise remain in \CERT \\
\midrule
New-regime creation & A candidate unmatched aggregate becomes a new regime only if it is separated from all stored shapes by~\eqref{eq:candidate-separated} and has an $\varepsilon_{\mathrm{opt}}$-certified arm. & \COMMIT after creation \\
\midrule
\COMMIT & Play the certified arm $A_t=a_j$ deterministically. Update the committed-arm detector~\eqref{eq:cusum}. Do not sample other arms and do not use committed-arm rewards to update the full shape. & Remain in \COMMIT unless~\eqref{eq:alarm} holds \\
\midrule
Alarm & If $G_t\geq h$, stop commitment, reset the detector, and return to balanced identification. & \ID \\
\bottomrule
\end{tabularx}
\caption{Library Identify-and-Commit with Thompson-Guided Certification. Thompson sampling is used only during certification; stable operation is deterministic.}
\label{tab:algorithm}
\end{table}

The same procedure can be written as explicit pseudocode.

\begin{center}
\fbox{%
\begin{minipage}{0.94\textwidth}
\textbf{Algorithm: Library Identify-and-Commit with Thompson-Guided Certification}
\begin{enumerate}[leftmargin=1.6em]
    \item Initialize the library $\Lcal=\varnothing$ and set the mode to \ID.
    \item In \ID, collect balanced blocks until a stable aggregate $C$ is available.
    \item Compute the match set $\mathcal M(C)$.
    \item If $\mathcal M(C)=\{j\}$, assign $C$ to regime $j$ and update both shape and certification statistics.
    \item If the matched regime has an $\varepsilon_{\mathrm{opt}}$-certified arm, enter \COMMIT. Otherwise enter \CERT.
    \item If $\mathcal M(C)=\varnothing$, treat $C$ as a candidate new regime and enter \CERT.
    \item In \CERT, sample arms by the top-two Thompson rule~\eqref{eq:ts-cert-action} from the posterior~\eqref{eq:working-posterior}, center each observed reward using the local level of $C$, and update the selected arm.
    \item Stop certification when condition~\eqref{eq:epsilon-cert} holds. For a candidate regime, also require separation from the library by~\eqref{eq:candidate-separated} before creating a new entry.
    \item In \COMMIT, play the certified arm at every round and update the detector $G_t$.
    \item If $G_t\geq h$, raise a drift alarm and return to \ID.
\end{enumerate}
\end{minipage}}
\end{center}

\section{Regret interpretation}

The correct theoretical statement is a high-probability regret decomposition. Let $G_T$ be the number of true stable episodes up to horizon $T$. Let $D_g$ be the delay between the $g$-th true drift and its detection. Let $T_{\mathrm{com}}$ be the number of rounds spent in \COMMIT. Let $\Delta_{\max}$ bound one-step regret.

On the event that the confidence intervals used for matching and certification are valid, any arm certified by~\eqref{eq:epsilon-cert} is at most $\varepsilon_{\mathrm{opt}}$ worse than the true best arm of the active regime. Therefore the regret satisfies
\begin{equation}
    R_T
    \leq
    R_{\mathrm{id}}(T)
    +
    R_{\mathrm{cert}}(T)
    +
    \Delta_{\max}\sum_{g=1}^{G_T}D_g
    +
    \varepsilon_{\mathrm{opt}}T_{\mathrm{com}}.
    \label{eq:regret}
\end{equation}
A high-probability statement adds the failure-event cost
\begin{equation}
    \delta T\Delta_{\max}.
\end{equation}
The term $R_{\mathrm{id}}(T)$ is the cost of balanced blocks used to identify regimes. The term $R_{\mathrm{cert}}(T)$ is the cost of certification samples. The certification sampling rule affects $R_{\mathrm{cert}}(T)$ by allocating certification pulls toward arms that are still plausibly optimal.

\begin{remark}[The guarantee does not depend on the certification sampling rule]
\label{rem:sampling-free}
The high-probability statement above---on the event that the bounds~\eqref{eq:cert-bounds} hold, every arm certified by~\eqref{eq:epsilon-cert} is at most $\varepsilon_{\mathrm{opt}}$-suboptimal, and hence~\eqref{eq:regret} holds---uses only the validity of the confidence bounds and the stopping rule~\eqref{eq:epsilon-cert}. It never uses how the certification pulls were allocated. Because each arm's reward fluctuations are conditionally sub-Gaussian (Assumption~\ref{ass:subgaussian}) and the bounds hold simultaneously for all arms and all sample counts through the union bound absorbed in the $8KT/\delta$ factor, they remain valid under adaptive, data-dependent sampling and under a data-dependent stopping time. Consequently the correctness guarantee and the decomposition~\eqref{eq:regret} hold verbatim for \emph{any} certification sampling rule, in particular for plain Thompson sampling ($\beta=1$) and for the top-two rule~\eqref{eq:ts-cert-action} with $\beta\in(0,1)$.
\end{remark}

\begin{remark}[Why a top-two rule]
\label{rem:toptwo}
The sampling rule affects only $R_{\mathrm{cert}}(T)$. The stopping rule~\eqref{eq:epsilon-cert} can fire only once $\max_{k\neq\widehat a}\UCB_k$ has fallen, which requires certification samples on the \emph{non-leading} arms. Plain Thompson sampling ($\beta=1$) allocates almost all pulls to the current leader and may leave these bounds wide, so certification can stall even when the leader is clearly best. The top-two rule with $\beta\in(0,1)$ spends a fixed fraction $1-\beta$ on a challenger, so every radius $B_{k}^{\mathrm{cert}}$ shrinks and~\eqref{eq:epsilon-cert} is met after a per-episode certification budget of order
\begin{equation}
    \sum_{k\in\A}
    \frac{\sigma^2\log(8KT/\delta)}
    {\big(\max\{\Delta_{g,k},\varepsilon_{\mathrm{opt}}\}\big)^2},
    \label{eq:cert-budget}
\end{equation}
the standard $\varepsilon_{\mathrm{opt}}$-optimal top-two allocation. This tightens $R_{\mathrm{cert}}(T)$ without changing any guarantee.
\end{remark}

If exact gaps are denoted by
\begin{equation}
    \Delta_{g,k}=\mu_{k^\star(g)}^{(g)}-\mu_k^{(g)},
\end{equation}
then the sample complexity of $\varepsilon_{\mathrm{opt}}$-certification scales with effective gaps of order
\begin{equation}
    \max\{\Delta_{g,k},\varepsilon_{\mathrm{opt}}\}.
\end{equation}
Thus nearly tied arms do not force the algorithm to continue sampling until exact separation is proved. The price is the explicit commitment term $\varepsilon_{\mathrm{opt}}T_{\mathrm{com}}$ in~\eqref{eq:regret}.

\section{Summary}

The algorithm separates the three statistical tasks. Balanced blocks identify regimes. Thompson-guided certification samples arms only while the best arm is uncertain. The stopping rule certifies practical optimality, not exact optimality. Once an arm is certified, the algorithm commits to it and uses a two-sided detector on the committed arm to detect drift. The only new conceptual parameter is $\varepsilon_{\mathrm{opt}}$, the smallest reward gap worth resolving. This modification prevents the algorithm from stalling in regimes where several arms have almost identical reward means.

\end{document}
